% Môn: Toán
% Lớp: 10
% Chương: Chương II. Bất phương trình và hệ bất phương trình bậc nhất hai ẩn
% Bài: Bài 3. Bất phương trình bậc nhất hai ẩn
% Số câu: 17
% App đọc file này trực tiếp — sửa rồi Commit trên GitHub, bấm Đồng bộ GitHub .tex

% Bài 3. Bất phương trình bậc nhất hai ẩn

% ===== Câu 1 =====
% ID: T10C2B3-01-DS
% Mức: TH
\dangbt{Nhận biết bất phương trình bậc nhất hai ẩn và nghiệm}
\begin{ex}
% ID: T10C2B3-01-DS
% Mức: TH
Cho bất phương trình bậc nhất hai ẩn $x-2y+2 \leq 0$.
\choiceTF
{\True $(1;4)$ là nghiệm của bất phương trình $x-2y+2 \leq 0$}
{$(0;3)$ không là nghiệm của bất phương trình $x-2y+2 \leq 0$}
{Đường thẳng $d \colon x-2y+2=0$ đi qua hai điểm $A(0;-1)$ và $B(2;0)$}
{\True Miền nghiệm của bất phương trình $x-2y+2 \leq 0$ là nửa mặt phẳng kể cả bờ $x-2y+2=0$, không chứa điểm $(0;0)$}
\loigiai{
\begin{itemchoice}
\item (Đúng) $(1;4)$ là nghiệm của bất phương trình $x-2y+2 \leq 0$. (Vì): Thay $(1;4)$ vào bất phương trình: $1 - 2(4) + 2 = 1 - 8 + 2 = -5$. Vì $-5 \leq 0$ nên $(1;4)$ là nghiệm của bất phương trình
\item (Sai) $(0;3)$ không là nghiệm của bất phương trình $x-2y+2 \leq 0$. (Vì): Thay $(0;3)$ vào bất phương trình: $0 - 2(3) + 2 = 0 - 6 + 2 = -4$. Vì $-4 \leq 0$ nên $(0;3)$ là nghiệm của bất phương trình
\item (Sai) Đường thẳng $d \colon x-2y+2=0$ đi qua hai điểm $A(0;-1)$ và $B(2;0)$. (Vì): Thay tọa độ điểm $A(0;-1)$ vào phương trình đường thẳng $d \colon x-2y+2=0$: $0 - 2(-1) + 2 = 2 + 2 = 4 \neq 0$. Điểm $A(0;-1)$ không thuộc đường thẳng $d$. Thay tọa độ điểm $B(2;0)$ vào phương trình đường thẳng $d \colon x-2y+2=0$: $2 - 2(0) + 2 = 2 + 2 = 4 \neq 0$. Điểm $B(2;0)$ không thuộc đường thẳng $d$
\item (Đúng) Miền nghiệm của bất phương trình $x-2y+2 \leq 0$ là nửa mặt phẳng kể cả bờ $x-2y+2=0$, không chứa điểm $(0;0)$. (Vì): ường thẳng $d \colon x-2y+2=0$ là bờ của nửa mặt phẳng. Thay tọa độ điểm $(0;0)$ vào biểu thức $x-2y+2$: $0 - 2(0) + 2 = 2$. Vì $2 \gt  0$ và bất phương trình là $x-2y+2 \leq 0$, nên miền nghiệm là nửa mặt phẳng không chứa điểm $(0;0)$ và bao gồm cả đường thẳng $d$
\end{itemchoice}
}
\end{ex}

% ===== Câu 2 =====
% ID: T10C2B3-01-TLN
% Mức: TH
\dangbt{Kiểm tra một cặp số là nghiệm của bất phương trình}
\begin{ex}
% ID: T10C2B3-01-TLN
% Mức: TH
Biết đường thẳng $d \colon 3x-4y-12=0$ đi qua hai điểm $A(0;a)$ và $B(b;0)$. Giá trị $a+b$ bằng bao nhiêu?
\shortans{$1$}
\loigiai{
Thay tọa độ $(0;a)$ vào $d$ ta được $3 \cdot 0 - 4\cdot a-12=0$ nên $a=-3$.
Thay tọa độ $(b;0)$ vào $d$ ta được $3 \cdot b - 4\cdot 0-12=0$ nên $b=4$.
Vậy $a+b=-3+4=1$.
}
\end{ex}

% ===== Câu 3 =====
% ID: T10C2B3-01-TN
% Mức: TH
\dangbt{Lập và giải bài toán thực tế bằng bất phương trình bậc nhất hai ẩn}
\begin{ex}
% ID: T10C2B3-01-TN
% Mức: TH
Cho bất phương trình $2x+3y-6 \leq 0\quad (1)$. Chọn khẳng định đúng trong các khẳng định sau
\choice
{Bất phương trình đã cho vô nghiệm.}
{Tập nghiệm của bất phương trình là đường thẳng $2x+3y-6=0$.}
{\True Bất phương trình đã cho có vô số nghiệm.}
{Tập nghiệm của bất phương trình là nửa mặt phẳng bờ $2x+3y-6=0$ không chứa gốc tọa độ.}
\loigiai{
Trên mặt phẳng tọa độ, đường thẳng $D.$.
Do đó, bất phương trình $(1)$ có vô số nghiệm
}
\end{ex}

% ===== Câu 4 =====
% ID: T10C2B3-02-DS
% Mức: TH
\dangbt{Kiểm tra một cặp số là nghiệm của bất phương trình}
\begin{ex}
% ID: T10C2B3-02-DS
% Mức: TH
An thích ăn hai loại trái cây là cam và xoài, mỗi tuần mẹ cho An $200\,000$ đồng để mua trái cây. Biết rằng giá cam là $15\,000$ đồng/kg, giá xoài là $30\,000$ đồng/kg. Gọi $x$, $y$ lần lượt là số kilôgam cam và xoài mà An có thể mua về sử dụng trong một tuần.
\choiceTF
{\True Trong tuần, số tiền An có thể mua cam là $15\,000x$, số tiền An có thể mua xoài là $30\,000 y$, ($x$, $y>0$)}
{Bất phương trình bậc nhất cho hai ẩn $x$, $y$ là $3 x+6 y \geq 40$}
{\True Cặp số $(5 ; 4)$ thỏa mãn bất phương trình bậc nhất cho hai ẩn $x$, $y$}
{An có thể mua $4$ kg cam, $5$ kg xoài trong tuần}
\loigiai{
\begin{itemchoice}
\item (Đúng) Trong tuần, số tiền An có thể mua cam là $15\,000x$, số tiền An có thể mua xoài là $30\,000 y$, ($x$, $y>0$). (Vì): Trong tuần, số tiền An có thể mua cam là $15\,000 x$, số tiền An có thể mua xoài là $30\,000 y$ ($x$, $y>0$
\item (Sai) Bất phương trình bậc nhất cho hai ẩn $x$, $y$ là $3 x+6 y \geq 40$. (Vì): Ta có bất phương trình $15\,000 x+30\,000 y \leq 200\,000 \Leftrightarrow 3 x+6 y \leq 40$. $\quad\ast)$
\item (Đúng) Cặp số $(5 ; 4)$ thỏa mãn bất phương trình bậc nhất cho hai ẩn $x$, $y$. (Vì): Xét $x=5$, $y=4$, thay vào bất phương trình $3\cdot5+6\cdot 4 \leq 40$ (đúng) nên $(5 ; 4)$ là một nghiệm của $\ast)$
\item (Sai) An có thể mua $4$ kg cam, $5$ kg xoài trong tuần. (Vì): Xét $x=4$, $y=5$, thay vào bất phương trình $3\cdot4+6\cdot 5 \leq 40$ (sai) nên $(4 ; 5)$ không là một nghiệm của $\ast)$
\end{itemchoice}
}
\end{ex}

% ===== Câu 5 =====
% ID: T10C2B3-02-TLN
% Mức: TH
\dangbt{Xác định đường biên và miền nghiệm}
\begin{ex}
% ID: T10C2B3-02-TLN
% Mức: TH
Bạn Nam tiết kiệm được $450$ nghìn đồng. Trong đợt ủng hộ các bạn học sinh đồng bào miền Trung bị lũ lụt vừa qua, bạn Nam đã ủng hộ $x$ tờ tiền loại $20$ nghìn đồng, $y$ tờ tiền loại $10$ nghìn đồng. Khi đó bất phương trình biểu diễn tổng số tiền mà bạn Nam đã ủng hộ có dạng $ax+y\le b$. Tính giá trị của biểu thức $P=2a-b$.
\shortans{$-41$}
\loigiai{
Tổng số tiền bạn Nam đã ủng hộ là $20x + 10y$ (nghìn đồng).
Số tiền bạn Nam có là $450$ nghìn đồng.
Bất phương trình biểu diễn tổng số tiền bạn Nam đã ủng hộ là $20x + 10y \le 450$.
Chia cả hai vế cho $10$, ta được $2x + y \le 45$.
So sánh với dạng $ax + y \le b$, ta có $a=2$ và $b=45$.
$P = 2a - b = 2(2) - 45 = 4 - 45 = -41$.
$-41$
}
\end{ex}

% ===== Câu 6 =====
% ID: T10C2B3-02-TN
% Mức: TH
\dangbt{Nhận biết bất phương trình bậc nhất hai ẩn và nghiệm}
\begin{ex}
% ID: T10C2B3-02-TN
% Mức: TH
Bất phương trình nào sau đây là bất phương trình bậc nhất hai ẩn?
\choiceTF
{$x^2+y^2 \leq 1$}
{$2x-y^3 \gt  0$}
{\True $x-y \leq 0$}
{$-\dfrac{2}{5} x-5^2 y \leq-\sqrt{15}$}
\loigiai{
\begin{itemchoice}
\item (Sai) $x^2+y^2 \leq 1$. (Vì): Bất phương trình $x^2+y^2 \leq 1$ chứa các số hạng bậc hai ($x^2$ và $y^2$), do đó không phải là bất phương trình bậc nhất hai ẩn
\item (Sai) $2x-y^3 \gt  0$. (Vì): Bất phương trình $2x-y^3 \gt  0$ chứa số hạng bậc ba ($y^3$), do đó không phải là bất phương trình bậc nhất hai ẩn
\item (Đúng) $x-y \leq 0$. (Vì): Bất phương trình $x-y \leq 0$ có dạng $ax+by \leq c$ với $a=1$, $b=-1$, $c=0$. Các hệ số $a$ và $b$ không đồng thời bằng 0, do đó đây là bất phương trình bậc nhất hai ẩn
\item (Sai) $-\dfrac{2}{5} x-5^2 y \leq-\sqrt{15}$.
\end{itemchoice}
}
\end{ex}

% ===== Câu 7 =====
% ID: T10C2B3-03-DS
% Mức: TH
\dangbt{Xác định đường biên và miền nghiệm}
\begin{ex}
% ID: T10C2B3-03-DS
% Mức: TH
Cho điểm $(-1 ; 2)$ và các bất phương trình $3 x-5 y<-15$, $2 x+y \le 0$, $3 x-9 y>7$, $-4 x+3 y \ge 5$.
\choiceTF
{\True $(-1 ; 2)$ không là một nghiệm của bất phương trình $3 x-5 y<-15$}
{\True $(-1 ; 2)$ là một nghiệm của bất phương trình $2 x+y \le 0$}
{$(-1 ; 2)$ là một nghiệm của bất phương trình $3 x-9 y>7$}
{\True $(-1 ; 2)$ là một nghiệm của bất phương trình $-4 x+3 y \ge 5$}
\loigiai{
\begin{itemchoice}
\item (Đúng) $(-1 ; 2)$ không là một nghiệm của bất phương trình $3 x-5 y<-15$. (Vì): Thay $(-1; 2)$ vào bất phương trình $3x - 5y \lt  -15$, ta được $3(-1) - 5(2) = -3 - 10 = -13$. Vì $-13 \not< -15$, nên $(-1; 2)$ không là nghiệm của bất phương trình $3x - 5y \lt  -15$
\item (Đúng) $(-1 ; 2)$ là một nghiệm của bất phương trình $2 x+y \le 0$. (Vì): Thay $(-1; 2)$ vào bất phương trình $2x + y \le 0$, ta được $2(-1) + 2 = -2 + 2 = 0$. Vì $0 \le 0$, nên $(-1; 2)$ là nghiệm của bất phương trình $2x + y \le 0$
\item (Sai) $(-1 ; 2)$ là một nghiệm của bất phương trình $3 x-9 y>7$. (Vì): Thay $(-1; 2)$ vào bất phương trình $3x - 9y \gt  7$, ta được $3(-1) - 9(2) = -3 - 18 = -21$. Vì $-21 \not> 7$, nên $(-1; 2)$ không là nghiệm của bất phương trình $3x - 9y \gt  7$
\item (Đúng) $(-1 ; 2)$ là một nghiệm của bất phương trình $-4 x+3 y \ge 5$. (Vì): Thay $(-1; 2)$ vào bất phương trình $-4x + 3y \ge 5$, ta được $-4(-1) + 3(2) = 4 + 6 = 10$. Vì $10 \ge 5$, nên $(-1; 2)$ là nghiệm của bất phương trình $-4x + 3y \ge 5$
\end{itemchoice}
}
\end{ex}

% ===== Câu 8 =====
% ID: T10C2B3-03-TLN
% Mức: TH
\begin{ex}
% ID: T10C2B3-03-TLN
% Mức: TH
Bạn Lan mang $150\,000$ đồng đi nhà sách để mua một số quyển tập và bút. Biết rằng giá một quyển tập là $8\,000$ đồng và giá của một cây bút là $6\,000$ đồng. Bạn Lan có thể mua được tối đa bao nhiêu quyển tập nếu bạn đã mua $10$ cây bút.
\shortans{$11$}
\loigiai{
Gọi $x$, $y$ lần lượt là số tập và số bút. Điều kiện $x>0$, $y>0$.
 
 
 • Bất phương trình biểu diễn số tập và bút có thể mua được phụ thuộc vào số tiền mang theo là $8\,000x+6\,000y \leq 150\,000$.
 
 • Bạn Lan có thể mua được tối đa số quyển tập nếu bạn đã mua $10$ cây bút là $8\,000x+6\,000 \cdot 10 \leq 150\,000 \Leftrightarrow x \leq 11,25$. 
 
 • Vì $x$ nguyên dương nên số quyển tập tối đa bạn Lan mua được $11$ quyển.
}
\end{ex}

% ===== Câu 9 =====
% ID: T10C2B3-03-TN
% Mức: TH
\dangbt{Xác định bất phương trình từ miền nghiệm}
\begin{ex}
% ID: T10C2B3-03-TN
% Mức: TH
Điểm $O(0 ; 0)$ thuộc miền nghiệm của bất phương trình nào dưới đây?
\choice
{$3x+3y \le -2$}
{$\dfrac{1}{3}x+2y+2 \le 0$}
{\True $2x+4y-2\lt 0$}
{$x+10y>0$}
\loigiai{
Thay $x=y=0$ vào bất phương trình $2x+4y-2\lt 0$, ta được $2\cdot 0+4 \cdot 0 -2\lt 0$ (đúng).
Vì vậy $O(0 ; 0)$ thuộc miền nghiệm.
}
\end{ex}

% ===== Câu 10 =====
% ID: T10C2B3-04-DS
% Mức: TH
\dangbt{Nhận biết bất phương trình bậc nhất hai ẩn và nghiệm}
\begin{ex}
% ID: T10C2B3-04-DS
% Mức: TH
Cho bất phương trình $3x+y>3 \qquad(1)$. Xét các mệnh đề sau:
\choiceTF
{Cặp số $(1;0)$ là nghiệm của bất phương trình $(1)$.}
{Cặp số $(0;0)$ là nghiệm của bất phương trình $(1)$.}
{Miền nghiệm của bất phương trình $(1)$ chứa gốc tọa độ $O(0;0)$.}
{Bất phương trình $3x+y>3$ luôn có miền nghiệm.}
\loigiai{
\begin{itemchoice}
\item (Sai) Cặp số $(1;0)$ là nghiệm của bất phương trình $(1)$. (Vì): Thay cặp số $(1;0)$ vào bất phương trình, ta có $3(1)+0 = 3$. Vì $3 \gt  3$ là sai nên cặp số $(1;0)$ không là nghiệm của bất phương trình
\item (Sai) Cặp số $(0;0)$ là nghiệm của bất phương trình $(1)$. (Vì): Thay cặp số $(0;0)$ vào bất phương trình, ta có $3(0)+0 = 0$. Vì $0 \gt  3$ là sai nên cặp số $(0;0)$ không là nghiệm của bất phương trình
\item (Sai) Miền nghiệm của bất phương trình $(1)$ chứa gốc tọa độ $O(0;0)$. (Vì): Từ kết quả của mệnh đề $B$, ta thấy gốc tọa độ $O(0;0)$ không là nghiệm của bất phương trình, do đó miền nghiệm của bất phương trình không chứa gốc tọa độ $O(0;0)$
\item (Đúng) Bất phương trình $3x+y>3$ luôn có miền nghiệm. (Vì): Bất phương trình bậc nhất hai ẩn $ax+by>c$ luôn có một miền nghiệm là một nửa mặt phẳng không kể bờ
\end{itemchoice}
}
\end{ex}

% ===== Câu 11 =====
% ID: T10C2B3-04-TN
% Mức: TH
\begin{ex}
% ID: T10C2B3-04-TN
% Mức: TH
Bất phương trình $3x-2(y-x+1)\gt 0$ tương đương với bất phương trình nào sau đây?
\choice
{$x-2y-2\gt 0$}
{\True $5x-2y-2\gt 0$}
{$5x-2y-1\gt 0$}
{$4x-2y-2\gt 0$}
\loigiai{
Ta có $3x-2(y-x+1)\gt 0\Leftrightarrow 3x-2y+2x-2\gt 0\Leftrightarrow 5x-2y-2\gt 0$.
}
\end{ex}

% ===== Câu 12 =====
% ID: T10C2B3-05-TN
% Mức: TH
\dangbt{Kiểm tra một cặp số là nghiệm của bất phương trình}
\begin{ex}
% ID: T10C2B3-05-TN
% Mức: TH
Cho bất phương trình $-2 x+3 y>3$. Khẳng định nào dưới đây $\textbf{sai}$?
\choice
{$(0 ; 0)$ không là nghiệm bất phương trình}
{\True $(-1 ; 1)$ không là nghiệm bất phương trình}
{$(0 ; 1)$ không là nghiệm bất phương trình}
{$(1 ; 3)$ là nghiệm bất phương trình}
\loigiai{
Thay $x=-1$, $y=1$ vào bất phương trình $2+3\gt 3$ (đúng).
Vì vậy $(-1 ; 1)$ là nghiệm bất phương trình.
}
\end{ex}

% ===== Câu 13 =====
% ID: T10C2B3-06-TN
% Mức: TH
\begin{ex}
% ID: T10C2B3-06-TN
% Mức: TH
Miền nghiệm của bất phương trình $3(x-1)+4(y-2)\lt 5x-3$ là nửa mặt phẳng chứa điểm
\choice
{$Q(-5;3)$}
{\True $O(0;0)$}
{$N(-4;2)$}
{$P(-2;2)$}
\loigiai{
Ta có $3(x-1)+4(y-2)\lt 5x-3\Leftrightarrow -x+2y-4\lt 0 \quad(\ast)$.
Thay $O(0;0)$ vào $(\ast)$ ta được $-0+1\cdot (-4)= -4 \lt  0$.
Do đó $O(0;0)$ thuộc miền nghiệm của bất phương trình.
}
\end{ex}

% ===== Câu 14 =====
% ID: T10C2B3-07-TN
% Mức: TH
\begin{ex}
% ID: T10C2B3-07-TN
% Mức: TH
Miền không gạch chéo (không kể bờ $d$) trong hình bên là miền nghiệm của bất phương trình nào trong các bất phương trình dưới đây?
\choice
{$2x+3y<6$}
{\True $2x+3y>6$}
{$\dfrac{x}{2}+\dfrac{y}{3}\gt 0$}
{$\dfrac{x}{2}+\dfrac{y}{3}\lt 1$}
\loigiai{
Xét đường thẳng $d\colon y=ax+b$ đi qua hai điểm $A(3,0)$ và $B(0,2)$.
Ta có $\left\{ \begin{aligned} 0=3a+b\ & 2=b \end{aligned} \right.\Leftrightarrow \left\{ \begin{aligned} a=-\dfrac{2}{3}\& b=2 \end{aligned} \right.\Rightarrow d\colon y=-\dfrac{2}{3}x+2\Leftrightarrow d\colon 2x+3y-6=0$.
Xét điểm $O(0;0)$, ta có $2\cdot 0+3\cdot 0-6=-6\lt 0$ nên điểm $O$ thuộc miền nghiệm của bất phương trình $2x+3y-6\lt 0$.
Vậy miền không gạch chéo (không kể bờ $d$) là miền nghiệm của bất phương trình $2x+3y-6\gt 0$ hay $2x+3y>6$.
}
\end{ex}

% ===== Câu 15 =====
% ID: T10C2B3-08-TN
% Mức: TH
\begin{ex}
% ID: T10C2B3-08-TN
% Mức: TH
Trong các bất phương trình sau đây, đâu là bất phương trình bậc nhất hai ẩn
\choice
{$2x^2-3x \geq 1$}
{\True $2x+y\leq 1$}
{$3x+\sqrt{y} \leq 0$}
{$3x+y=1$}
\loigiai{
Theo định nghĩa $2x+y\leq 1$ là bất phương trình bậc nhất hai ẩn.
}
\end{ex}

% ===== Câu 16 =====
% ID: T10C2B3-09-TN
% Mức: TH
\begin{ex}
% ID: T10C2B3-09-TN
% Mức: TH
Trong các bất phương trình sau, bất phương trình nào là bất phương trình bậc nhất hai ẩn?
\choice
{$2x^2+3y>4$}
{$xy+x<5$}
{\True $3^2x+4^3y \geq 6$}
{$x+y^3 \leq 3$}
\loigiai{
Bất phương trình $3^2x+4^3y \geq 6$ là bất phương trình bậc nhất hai ẩn.
}
\end{ex}

% ===== Câu 17 =====
% ID: T10C2B3-10-TN
% Mức: TH
\begin{ex}
% ID: T10C2B3-10-TN
% Mức: TH
Cho bất phương trình bậc nhất hai ẩn $2 x-5 y \leq 8$. Khẳng định nào dưới đây là khẳng định $\textbf{sai}$?
\choice
{$(3 ;-4)$ là một nghiệm của bất phương trình}
{\True $(-2 ; 2)$ không là một nghiệm của bất phương trình}
{$(-3 ;-1)$ là một nghiệm của bất phương trình}
{$(5 ; 0)$ không là một nghiệm của bất phương trình}
\loigiai{
Thay $(-2 ; 2)$ vào bất phương trình $2 x-5 y \leq 8$ ta được $2 \cdot (-2)-5\cdot 2 \leq 8 \Leftrightarrow-14 \leq 8$
  (đúng), nên $(-2 ; 2)$ là một nghiệm của bất phương trình $2 x-5 y \leq 8$.
}
\end{ex}
